\section{Model}\label{sec:model}

\subsection{Users, integration, and future usage}

A monopoly provider faces two classes of potential enterprise task integrations, $j\in\{L,H\}$. Each class has unit potential mass. A task of class $j$ draws a sunk integration cost $k$ independently from a uniform distribution on $[0,K_j]$, with $K_j>0$. If it integrates, it receives a current benefit $b_j>0$ before any future repricing. This benefit is realized before the provider's later tariff decision and cannot be expropriated by that decision.

Future gross value from usage $q\geq0$ is
\begin{equation}
 v_j(q)=a_jq-\frac{q^2}{2},
 \label{eq:utility}
\end{equation}
where
\begin{equation}
 0<c<a_L<a_H,
 \qquad d\equiv a_H-a_L>0.
 \label{eq:primitive-order}
\end{equation}
The provider incurs real inference cost $cq$. The two classes therefore differ in their future value of AI usage, with class $H$ the high-use class.

The provider posts a common two-part tariff
\begin{equation}
 T(q)=F+pq.
 \label{eq:tariff}
\end{equation}
The \emph{flat architecture} sets $p=0$. The \emph{metered/hybrid architecture} sets $p>0$ and incurs a real activation cost $\mu\geq0$. Within either architecture the provider chooses the feasible tariff parameters globally. The baseline does not allow discriminatory menus across classes, arbitrary nonlinear tariffs, competition, congestion, capacity choice, or an intermediary layer.

Conditional on participation under tariff $(F,p)$, type $j$ chooses
\begin{equation}
 q_j(p)=\max\{a_j-p,0\},
 \label{eq:demand}
\end{equation}
and its indirect gross surplus before the fixed fee is
\begin{equation}
 S_j(p)=\frac{(a_j-p)^2}{2}
 \label{eq:indirect-surplus}
\end{equation}
whenever $p<a_j$. Participation is voluntary.

\subsection{Timing and equilibrium concept}

The timing is:
\begin{enumerate}
 \item potential tasks simultaneously decide whether to incur their sunk integration costs;
 \item integrated tasks receive the current benefits $b_j$;
 \item the provider observes the installed masses $(l,h)$;
 \item the provider chooses the future tariff architecture and then the tariff parameters permitted by that architecture;
 \item integrated tasks choose future participation and usage.
\end{enumerate}

Users are forward-looking and atomless. A single task's integration decision does not change the aggregate installed state or the provider's continuation choice. Equilibrium therefore combines rational expectations at the integration stage with subgame-perfect continuation optimization by the provider. If the provider mixes between architectures, users take the aggregate mixing probability as given when deciding whether to integrate.

\subsection{The regular both-served region}

The headline results are not global over all positive parameters. They are stated on a certified regular region $\mathcal R$ in which the economically relevant continuation branch serves both classes under the globally optimal tariff within each architecture. Define the flat-price rent difference
\begin{equation}
 R_F\equiv S_H(0)-S_L(0)=\frac{a_H^2-a_L^2}{2}.
 \label{eq:RF}
\end{equation}
The region $\mathcal R$ imposes the following restrictions:
\begin{align}
 &K_L,K_H>0, \qquad l\equiv \frac{b_L}{K_L}\in(0,1),
 \label{eq:R1}\\
 &dc<b_H+R_F<K_H,
 \label{eq:R2}
\end{align}
together with two substantive continuation conditions: the relevant interior metered price remains below $a_L$, and on the full installed-state interval used by the equilibrium construction the globally optimal flat and metered continuations both serve $L$ and $H$ rather than switching to an $H$-only or no-service active set.

The last condition is essential. It permits closed-form analysis of the regular branch but is not innocuous. Section~\ref{sec:robustness} retains an exact outside-$\mathcal R$ counterexample showing why the global-all-parameters formulation is false. Equality cases at the architecture thresholds are also excluded from the strict headline theorem because the provider is indifferent there.
